\documentclass[10pt, letterpaper]{article}

% Packages:
\usepackage[
    ignoreheadfoot, % set margins without considering header and footer
    top=2 cm, % seperation between body and page edge from the top
    bottom=2 cm, % seperation between body and page edge from the bottom
    left=2 cm, % seperation between body and page edge from the left
    right=2 cm, % seperation between body and page edge from the right
    footskip=1.0 cm, % seperation between body and footer
    % showframe % for debugging 
]{geometry} % for adjusting page geometry
\usepackage{titlesec} % for customizing section titles
\usepackage{tabularx} % for making tables with fixed width columns
\usepackage{array} % tabularx requires this
\usepackage[dvipsnames]{xcolor} % for coloring text
\definecolor{primaryColor}{RGB}{0, 79, 144} % define primary color
\usepackage{enumitem} % for customizing lists
\usepackage{fontawesome5} % for using icons
\usepackage{amsmath} % for math
\usepackage[
    pdftitle={Wei-Hsuan Cheng's CV},
    pdfauthor={Wei-Hsuan Cheng},
    pdfcreator={LaTeX with RenderCV},
    colorlinks=true,
    urlcolor=primaryColor
]{hyperref} % for links, metadata and bookmarks
\usepackage[pscoord]{eso-pic} % for floating text on the page
\usepackage{calc} % for calculating lengths
\usepackage{bookmark} % for bookmarks
\usepackage{lastpage} % for getting the total number of pages
\usepackage{changepage} % for one column entries (adjustwidth environment)
\usepackage{paracol} % for two and three column entries
\usepackage{ifthen} % for conditional statements
\usepackage{needspace} % for avoiding page brake right after the section title
\usepackage{iftex} % check if engine is pdflatex, xetex or luatex

% Ensure that generate pdf is machine readable/ATS parsable:
\ifPDFTeX
    \input{glyphtounicode}
    \pdfgentounicode=1
    % \usepackage[T1]{fontenc} % this breaks sb2nov
    \usepackage[utf8]{inputenc}
    \usepackage{lmodern}
\fi



% Some settings:
\AtBeginEnvironment{adjustwidth}{\partopsep0pt} % remove space before adjustwidth environment
\pagestyle{empty} % no header or footer
\setcounter{secnumdepth}{0} % no section numbering
\setlength{\parindent}{0pt} % no indentation
\setlength{\topskip}{0pt} % no top skip
\setlength{\columnsep}{0cm} % set column seperation
\makeatletter
\let\ps@customFooterStyle\ps@plain % Copy the plain style to customFooterStyle
\patchcmd{\ps@customFooterStyle}{\thepage}{
    \color{gray}\textit{\small Wei-Hsuan Cheng - Page \thepage{} of \pageref*{LastPage}}
}{}{} % replace number by desired string
\makeatother
\pagestyle{customFooterStyle}

\titleformat{\section}{\needspace{4\baselineskip}\bfseries\large}{}{0pt}{}[\vspace{1pt}\titlerule]

\titlespacing{\section}{
    % left space:
    -1pt
}{
    % top space:
    0.3 cm
}{
    % bottom space:
    0.2 cm
} % section title spacing

\renewcommand\labelitemi{$\circ$} % custom bullet points
\newenvironment{highlights}{
    \begin{itemize}[
        topsep=0.10 cm,
        parsep=0.10 cm,
        partopsep=0pt,
        itemsep=0pt,
        leftmargin=0.4 cm + 10pt
    ]
}{
    \end{itemize}
} % new environment for highlights

\newenvironment{highlightsforbulletentries}{
    \begin{itemize}[
        topsep=0.10 cm,
        parsep=0.10 cm,
        partopsep=0pt,
        itemsep=0pt,
        leftmargin=10pt
    ]
}{
    \end{itemize}
} % new environment for highlights for bullet entries


\newenvironment{onecolentry}{
    \begin{adjustwidth}{
        0.2 cm + 0.00001 cm
    }{
        0.2 cm + 0.00001 cm
    }
}{
    \end{adjustwidth}
} % new environment for one column entries

\newenvironment{twocolentry}[2][]{
    \onecolentry
    \def\secondColumn{#2}
    \setcolumnwidth{\fill, 4.5 cm}
    \begin{paracol}{2}
}{
    \switchcolumn \raggedleft \secondColumn
    \end{paracol}
    \endonecolentry
} % new environment for two column entries

\newenvironment{header}{
    \setlength{\topsep}{0pt}\par\kern\topsep\centering\linespread{1.5}
}{
    \par\kern\topsep
} % new environment for the header

\newcommand{\placelastupdatedtext}{% \placetextbox{<horizontal pos>}{<vertical pos>}{<stuff>}
  \AddToShipoutPictureFG*{% Add <stuff> to current page foreground
    \put(
        \LenToUnit{\paperwidth-2 cm-0.2 cm+0.05cm},
        \LenToUnit{\paperheight-1.0 cm}
    ){\vtop{{\null}\makebox[0pt][c]{
        \small\color{gray}\textit{Last updated in Feb 2025}\hspace{\widthof{Last updated in Feb 2025}}
    }}}%
  }%
}%

% save the original href command in a new command:
\let\hrefWithoutArrow\href

% new command for external links:
\renewcommand{\href}[2]{\hrefWithoutArrow{#1}{\ifthenelse{\equal{#2}{}}{ }{#2 }\raisebox{.15ex}{\footnotesize \faExternalLink*}}}


\begin{document}
    \newcommand{\AND}{\unskip
        \cleaders\copy\ANDbox\hskip\wd\ANDbox
        \ignorespaces
    }
    \newsavebox\ANDbox
    \sbox\ANDbox{}

    \placelastupdatedtext
    \begin{header}
        \textbf{\fontsize{24 pt}{24 pt}\selectfont Wei-Hsuan Cheng}

        \vspace{0.3 cm}

        \normalsize


%%%%%%%%%% Contact info %%%%%%%%%%
        % Location
        \mbox{{\color{black}\footnotesize\faMapMarker*}\hspace*{0.13cm}Taipei, Taiwan}%
        \kern 0.25 cm%
        \AND%
        \kern 0.25 cm%
        % Email
        \mbox{\hrefWithoutArrow{mailto:johnathancheng0125@gmail.com}{\color{black}{\footnotesize\faEnvelope[regular]}\hspace*{0.13cm}johnathancheng0125@gmail.com}/\hrefWithoutArrow{mailto:r11631045@ntu.edu.tw}{\color{black}{\footnotesize\faEnvelope[regular]}\hspace*{0.13cm}r11631045@ntu.edu.tw}}%
        \kern 0.25 cm%
        \AND%
        \kern 0.25 cm%
        % TEL
        % \mbox{\hrefWithoutArrow{tel:+90-541-999-99-99}{\color{black}{\footnotesize\faPhone*}\hspace*{0.13cm}0541 999 99 99}}%
        % \kern 0.25 cm%
        % \AND%
        % \kern 0.25 cm%
        % LinkedIn
        \mbox{\hrefWithoutArrow{https://www.linkedin.com/in/wei-hsuan-cheng-582486264/}{\color{black}{\footnotesize\faLinkedinIn}\hspace*{0.13cm}Wei-Hsuan Cheng}}%
        \kern 0.25 cm%
        \AND%
        \kern 0.25 cm%
        % GitHub account
        \mbox{\hrefWithoutArrow{https://github.com/wei-hsuan-cheng}{\color{black}{\footnotesize\faGithub}\hspace*{0.13cm}wei-hsuan-cheng}}%
        \kern 0.25 cm%
        \AND%
        \kern 0.25 cm%
        % % Personal website
        % \mbox{\hrefWithoutArrow{https://wei-hsuan-cheng.github.io}{\color{black}{\footnotesize\faGithub}\hspace*{0.13cm}wei-hsuan-cheng.github.io}}%
    \end{header}

    \vspace{0.3 cm - 0.3 cm}


%%%%%%%%%% Education %%%%%%%%%%
    \section{Education}
        
        \begin{twocolentry}{
            
            
        \textit{(Expected) Spring 2025}}
            \textbf{National Taiwan University (NTU)}

            \textit{Master in Biomechatronics Engineering}
        \end{twocolentry}

        \vspace{0.10 cm}
        \begin{onecolentry}
            \begin{highlights}
                \item Overall GPA: 4.12/4.3
                \item Scholarships: 2023 Master Student Scholarship (First Place) from Dept. Biomechatronics Engineering, NTU
                \item \textbf{Coursework:} Robotics, Optimal Control, Stochastic Control, Robust Control, System Identification, Linear Systems
                \item \textbf{Master Thesis:} Application of Vision- and Force-Based Supervisory Teleoperation for Robotic Dexterous Manipulation Tasks in Industrial Surveillance
            \end{highlights}
        \end{onecolentry}

        \vspace{0.2 cm}

        \begin{twocolentry}{
            
        \textit{June 2022}}
            \textbf{National Tsing Hua University (NTHU)}
            
            \textit{Bachelor in Interdisciplinary Programme of Nuclear Science (IPNS)}
        \end{twocolentry}
        \vspace{0.10 cm}
        \begin{onecolentry}
            \begin{highlights}
                \item Specialising in Energy Engineering and Power Mechanical Engineering.
            \end{highlights}
        \end{onecolentry}


%%%%%%%%%% Skills %%%%%%%%%%
    \section{Skills}

        \begin{onecolentry}
            \textbf{Programming:} C++, Python, MATLAB, JavaScript, LaTeX
        \end{onecolentry}
        \vspace{0.2 cm}
        \begin{onecolentry}
            \textbf{Robotics Research Interests:} Motion Planning and Control, Kalman Filtering, Computer Vision, Vision- and Force-based Control, Application of Conformal Geometric Algebra (CGA) in Robotics
        \end{onecolentry}
        \vspace{0.2 cm}
        \begin{onecolentry}
            \textbf{Software/Hardware/Tools:} ROS 1 \& 2, MoveIt2!, PCL, PyQt5, Git, Docker, SolidWorks, Fusion 360, Nvidia Jetson, STM32 MCU 
        \end{onecolentry}


%%%%%%%%%% Research Experience %%%%%%%%%%
    \section{Research Experience}
        
        \begin{twocolentry}{
        \textit{Taipei, Taiwan}    
            
        \textit{Dec 2023 – Present}}
            \textbf{\href{https://dnc.ntu.edu.tw/en}{Delta-NTU Joint Research \& Development Centre}}
            
            \textit{Project leader in the collaboration research project of NTU and the CTO Office in Delta Electronics, Inc.}
        \end{twocolentry}

        \vspace{0.10 cm}
        \begin{onecolentry}
        \textbf{Supervisory Teleoperation of Manipulator for Surveillance}
            \begin{highlights}
                \item Develop a supervisory tele-manipulation system for industrial surveillance, with hierarchical-autonomy-level structure (teleop-supervised-autonomous). A human-robot interface (HRI) is designed for the operator to teleoperate remote robot via wireless communication and perform dexterous manipulation tasks in industrial surveillance scenario (\textit{e.g.} buttons, switches, and valves operations).
                \item Apply both vision- and force-based control strategies for robotic arm dexterous manipulation tasks (preparing for master thesis and submitting to \textbf{RA-L} journal).
            \end{highlights}
        \end{onecolentry}


        \vspace{0.2 cm}

        \begin{twocolentry}{
        \textit{Taipei, Taiwan}    
            
        \textit{Sept 2022 – Present}}
            \textbf{\href{http://rmml.bime.ntu.edu.tw/nturmmle.html}{Robots and Medical Mechatronics Lab (RMML)}}
            
            \textit{Master student researcher, advised by \hrefWithoutArrow{http://rmml.bime.ntu.edu.tw/profe.html}{Prof. Ping-Lang Yen}}
            
        \end{twocolentry}

        \vspace{0.10 cm}
        \begin{onecolentry}
        \textbf{Project Experiences and Contributions}
            \begin{highlights}
                \item Agricultural Robot:
                \begin{itemize}
                    \item Improve EKF-SLAM for agricultural tracked mobile robot by involving IMU bias estimation and on-line state covariance update.
                    \item Apply vision-based control scheme and object 6D pose estimation for robotic tomato harvesting.
                \end{itemize}
                \item Surgical Robot:
                \begin{itemize}
                    \item Improve attitude and angular velocity estimation, and the visual servo control scheme for surgical robot end-effector using multi-sensor fusion and a quaternion-based EKF.
                \end{itemize}
            \end{highlights}
        \end{onecolentry}



%%%%%%%%%% Publications %%%%%%%%%%
    % \section{Publications}

    %     \begin{samepage}
    %         \begin{twocolentry}{
    %             Jan 2004
    %         }
    %             \textbf{3D Finite Element Analysis of No-Insulation Coils}

    %             \vspace{0.10 cm}

    %             \mbox{Frodo Baggins}, \mbox{\textbf{\textit{Wei-Hsuan Cheng}}}, \mbox{Samwise Gamgee}
    %         \end{twocolentry}


    %         \vspace{0.10 cm}

    %         \begin{onecolentry}
    %     \href{https://doi.org/10.1109/TASC.2023.3340648}{10.1109/TASC.2023.3340648}
    %         \end{onecolentry}
    %     \end{samepage}


%%%%%%%%%% Selected Projects %%%%%%%%%% 
    \section{Selected Projects}

        \begin{twocolentry}{
        \textit{May 2023}}
            \textbf{Attitude and Angular Velocity Estimation of Surgical Robot Using Sensor Fusion with Optical Tracker, IMU, and Quaternion-based Extended Kalman Filter}
        \end{twocolentry}
        \vspace{0.10 cm}
        \begin{onecolentry}
            \begin{highlights}
                \item A multi-sensor fusion framework is used for estimating the attitude and angular velocity of surgical robot end-effector, for better robustness (against occlusions) and tracking precision in minimally invasive surgeries.
                \item An quaternion-based EKF algorithm is derived based on a constant angular velocity model.
                \item The EKF algorithm is simulated in MATLAB and implemented in ROS 2 C++.
            \end{highlights}
        \end{onecolentry}

        \vspace{0.2 cm}

        \begin{twocolentry}{
        \hrefWithoutArrow{https://github.com/wei-hsuan-cheng/kendo_robot}{{\footnotesize\faGithub}\hspace*{0.13cm}kendo\_robot}\textit{, Dec 2023}
        }
            \textbf{Kendo Robot- Combining Robotics Technologies with Martial Art}
        \end{twocolentry}
        \vspace{0.10 cm}
        \begin{onecolentry}
            \begin{highlights}
                \item A supervisory teleoperation kendo robot system. The robotic arm holding kendo sword tracks the opponent’s pose and wait for the operator’s command to attack on different body positions.
                \item Detect and track the 3D human body by DL-based real-time human pose estimation algorithm and RGB-D alignment. An interactive GUI is built using JavaScript for real-time visualisation and robotic motion planning.
                \item Robotic arm IK solver and sword motion planner are designed based on a new mathematical framework, conformal geometric algebra (CGA).
                \item Developed using ROS 1 Python, and a web-based interacture GUI uisng JavaScript. ROS-JavaScript communication through web socket.
            \end{highlights}
        \end{onecolentry}

        \vspace{0.2 cm}

        \begin{twocolentry}{
        \hrefWithoutArrow{https://github.com/wei-hsuan-cheng/cga_ik}{{\footnotesize\faGithub}\hspace*{0.13cm}cga\_ik}\textit{, Nov 2023}
        }
            \textbf{CGA IK- Closed-Form Inverse Kinematics Solver based on CGA}
        \end{twocolentry}
        \vspace{0.10 cm}
        \begin{onecolentry}
            \begin{highlights}
                \item Aim to promote a new mathematical framework \textbf{conformal geometric algebra (CGA)} for robotics applications. 
                \item Implement a closed-form IK solver for 6-DoF collaborative robots based on CGA.
                \item Simulated and visualised in a web-based interactive GUI using JavaScript. Implemented in  ROS 2 Python and tested on a real robot.
            \end{highlights}
        \end{onecolentry}

        \vspace{0.2 cm}

        \begin{twocolentry}{
        \textit{Dec 2022}}
            \textbf{Furuta Pendulum Balancing Control and Disturbance Rejection via State-Observer Feedback and Internal Model Design}
        \end{twocolentry}
        \vspace{0.10 cm}
        \begin{onecolentry}
            \begin{highlights}
                \item Stabilise the Furuta pendulum system by the design of state-observer feedback controller. An internal model is ustilised to reject the sinusoidal process noise in the input signal of maxon motor. 
                \item Course project of EE5128 Linear Systems @ NTU EE. System dynamics and control algorithm are simulated in Simulink, Simscape, and MATLAB.
            \end{highlights}
        \end{onecolentry}

%%%%%%%%%% Teaching Assistants %%%%%%%%%%
    \section{Teaching Assistants (TAs)}
    
        \begin{twocolentry}{
        \textit{Autumn 2022}}
            \textbf{Engineering Mathematics (I)}
        \end{twocolentry}

        \vspace{0.2 cm}

        \begin{twocolentry}{
        \textit{Autumn 2022}}
            \textbf{Introduction to Foreign Literature on Bio-Systems Engineering}
        \end{twocolentry}
        
%%%%%%%%%% Awards %%%%%%%%%%
    \section{Awards}
    
        \begin{twocolentry}{
        \textit{Dec 2023}}
            \textbf{Robotics Competition Winner (Kendo Robot)} as a team leader in the course project of CSIE5074 Robotics, offered by Prof. Li-Chen Fu @ NTU EE.
        \end{twocolentry}


\end{document}